% Source: Wald GR, physical PDF p. 255
%         (printed p. 242).
% Coverage: Chapter 9 problems 1--3.
% Figures/tables/footnotes: no figures, tables, or footnotes.
% Status: source-reviewed and compile-clean.
% Uncertainties: none.

\chapterproblems

\begin{enumerate}[label=\arabic*.,leftmargin=*]
\item Let \((M,g_{ab})\) be a spacetime with an everywhere timelike Killing field
\(\xi^a\). Let \(\mu\) be an arbitrary timelike curve parameterized by proper time,
\(\tau\), and let \(u^a\) denote the unit tangent to \(\mu\). Define
\(E=-u_a\xi^a\).
  \begin{enumerate}[label=\alph*),leftmargin=2em]
  \item Show that \(\lvert\dd E/\dd\tau\rvert\leq aE\), where \(a\) denotes the
  magnitude of the acceleration of \(\mu\).

  \item Use the result of part (a) to show that for a timelike curve with bounded
  integrated acceleration, i.e., \(\int a\,\dd\tau<\infty\), \(E\) can change only by
  a finite amount. Use this result to prove that \(\lvert\xi_a\xi^a\rvert\) is bounded
  along every timelike curve of bounded integrated acceleration.
  \end{enumerate}

The condition \(\int a\,\dd\tau<\infty\) is satisfied by all world lines whose
acceleration is produced by a physically realizable rocket ship. Thus, part (b) shows
that even if observers are equipped with rocket ships, they cannot reach spacetime
singularities where the norm of the timelike Killing field diverges to \(-\infty\), as
occurs, for example, in the negative mass Schwarzschild solution; see Chakrabarti,
Geroch, and Liang (1983).

\item Let \(M\) be the torus \((S^1\times S^1)\) and define the Lorentz metric
\(g_{ab}\) by (Misner 1963)
\[
  \dd s^2=\cos x(\dd y^2-\dd x^2)+2\sin x\,\dd x\dd y ,
\]
where the angular coordinates \(x,y\) have ranges \(0\leq x\leq2\pi\),
\(0\leq y\leq2\pi\). Show that the closed curves defined by \(x=\pi/2\) and
\(x=3\pi/2\) are null geodesics with affine parameter, \(\lambda\), related to the
\(y\)-coordinate by \(\tfrac12\lambda\,\dd y/\dd\lambda=1\). Hence, show that if
one traverses either of these curves infinitely many times in the negative
\(y\)-direction, only a finite amount of affine parameter is used up. Note that when
one parallelly transports the tangent to the geodesic around one cycle in the negative
\(y\)-direction, the tangent comes back larger by the factor of \(e^\pi\). Thus, this
compact manifold with smooth metric is null geodesically incomplete. Timelike and
spacelike geodesics similarly wind around the torus infinitely many times in the
\(y\)-direction using up only a finite amount of affine parameter as they approach
\(x=\pi/2\) and \(x=3\pi/2\), so \((M,g_{ab})\) also is timelike and spacelike
geodesically incomplete.

\item Eliminate the hypothesis of strong causality in Theorem~9.5.2 as follows.
Under the assumption that all past directed timelike geodesics have length greater
than \(3/\lvert C\rvert\), it was shown that (a) for \(p\in H^-(S)\) there exists a
maximum length geodesic, \(\gamma_p\), from \(p\) to \(S\) and (b) \(H^-(S)\) is
compact. Define \(T:H^-(S)\longrightarrow\mathbb R\) by
\(T(p)=\tau[\gamma_p]\). (i) Show that \(T\) is continuous and, hence, achieves a
minimum value on \(H^-(S)\). (ii) Show that \(T\) strictly decreases along each
future directed null geodesic generator of \(H^-(S)\). Thus, obtain a contradiction
with (i).
\end{enumerate}

\setcounter{footnote}{\chapterninepreviousfootnote}
\endgroup
